\usepackage{epigraph}
\setlength{\epigraphwidth}{0.7\textwidth}
\setlength{\epigraphrule}{0pt}

\usepackage{hyperref}
\hypersetup{
  colorlinks=true,
  urlcolor=blue,
  linkcolor=black,
  plainpages=false,
  pageanchor=false,
}
\usepackage{xurl}

\usepackage[hang,flushmargin]{footmisc} % 可选：让脚注悬挂缩进、贴边
\renewcommand{\thefootnote}{[\arabic{footnote}]}

% 控制引文文字的对齐方式（这才是你需要的）
\renewcommand{\textflush}{flushright}

% 控制出处文字的对齐方式
\renewcommand{\sourceflush}{flushright}

% ===== 编号深度 =====
\setcounter{secnumdepth}{3}

% ===== 字体 =====
% \usepackage{fontspec}
% \usepackage{xeCJK}

% \setmainfont{Sarasa Fixed SC}[
%     BoldFont = Sarasa Fixed SC Bold
% ]
% \setCJKmainfont{Sarasa Fixed SC}[
%     BoldFont = Sarasa Fixed SC Bold
% ]

% ===== 页边距 =====
\usepackage{geometry}
\geometry{a4paper, left=2.5cm, right=2.5cm, top=3cm, bottom=2.5cm}

% ===== 章标题格式 =====
\usepackage{titlesec}
\titleformat{\chapter}[block]
  {\raggedright\Huge\bfseries}
  {第\thechapter 章\quad}
  {0pt}
  {}

\titleformat{\section}[block]
  {\raggedright\Large\bfseries}
  {\thesection\quad}
  {0pt}
  {}

% ===== 页眉页脚 =====
\usepackage{fancyhdr}
\setlength{\headheight}{13pt}

% 让 subsection 写入 \rightmark
\makeatletter
\renewcommand{\subsectionmark}[1]{%
  \markright{\thesubsection\quad #1}%
}
\makeatother

% 正文页眉样式
\fancypagestyle{fancy}{%
  \fancyhf{}%
  \fancyhead[LE,LO]{\rightmark}%
  \fancyhead[RE,RO]{\thepage}%
  \renewcommand{\headrulewidth}{0pt}%
}

% 目录页样式：无页眉，保留页码（如需无页码，删掉 fancyfoot 那行）
\fancypagestyle{tocstyle}{%
  \fancyhf{}%
  \fancyfoot[C]{\thepage}%
  \renewcommand{\headrulewidth}{0pt}%
}

% 全局默认用 fancy，但目录处会显式切到 tocstyle
\pagestyle{fancy}  % 去掉引文和正文之间的横线

% 版权声明的年限
\newif\ifsameyear
\newcommand{\copyrightyear}{%
  \ifnum\year=2026\relax
    2026%
  \else
    2026--\the\year
  \fi
}

\newcommand{\todo}[1]{\textbf{[TODO: #1]}}
\newcommand{\draft}[1]{\textit{#1}}